\documentclass[../main.tex]{subfiles}

\begin{document}

\ifSubfilesClassLoaded{

    %\tableofcontents
    
    %\setcounter{chapter}{}
}{}

\chapter{ De Rham Cohomology }
% 代靖涵 25120222201319

\section{The De Rham Cohomology Groups}
\begin{definition}{}{}
    Let $M$ be a smooth manifold with or without boundary, and let $p$ be a nonnegative integer. Because $d: \Omega^{p}(M) \rightarrow \Omega^{p+1}(M)$ is linear, its kernel and image are linear subspaces. We define
\[
\begin{aligned}
\mathcal{Z}^{p}(M) &= \operatorname{Ker}\left(d: \Omega^{p}(M) \rightarrow \Omega^{p+1}(M)\right)=\{\text {closed } p \text {-forms on } M\}, \\
\mathcal{B}^{p}(M) &= \operatorname{Im}\left(d: \Omega^{p-1}(M) \rightarrow \Omega^{p}(M)\right)=\{\text {exact } p \text {-forms on } M\} .
\end{aligned}
\]
By convention, we consider $\Omega^{p}(M)$ to be the zero vector space when $p<0$ or $p>n=\operatorname{dim} M$, so that, for example, $\mathcal{B}^{0}(M)=0$ and $\mathcal{Z}^{n}(M)=\Omega^{n}(M)$.
\end{definition}

\begin{definition}{}{}
    The fact that every exact form is closed implies that $\mathcal{B}^p(M) \subseteq \mathcal{Z}^p(M)$. Thus, it makes sense to define the \dfntxt{de Rham cohomology group in degree $p$} (or the \dfntxt{$p$th de Rham group}) \dfntxt{of $M$} to be the quotient vector space
\[
H_{\mathrm{dR}}^p(M) = \frac{\mathcal{Z}^p(M)}{\mathcal{B}^p(M)}.
\]
\end{definition}
\end{document}